\section{引言}

在离散系统中，“信息守恒 + 局域规则 + 可视化验证”三者若能同时成立，往往意味着该系统可作为更复杂构造的可复用核。
本文关注一个极简但封闭的情形：仅使用 Zeckendorf 表示（无相邻 $1$ 的 Fibonacci 数字系统）作为稳定态约束，定义两种基本操作——\textbf{展开}与\textbf{折叠}——并要求其在有限分辨率 $m=6$ 下闭合。

Zeckendorf 定理保证每个正整数都存在且唯一地表示为不相邻 Fibonacci 数之和，从而给出一个天然的正规形与归一化算子 \cite{WikipediaZeckendorf}.
Fenwick 从局域校正模式角度系统性总结了 Zeckendorf 算术所需的有限局域重写规则族 \cite{Fenwick2003ZeckendorfArithmetic}；
Ahlbach--Usatine--Frougny--Pippenger 则证明 Zeckendorf 加法可在线性时间内完成，并可由线性规模、对数深度的组合逻辑网络实现 \cite{AhlbachEtAl2013ZeckendorfArithmetic}，
为本文的“固定相位局域网络/元胞自动机”编译提供了标准参考。

\paragraph{本文贡献}
本文贡献可概括为：
\begin{itemize}
  \item (C1) 给出 $m=6$ 下 $\M_6\to\X_6$ 的有限纤维化，并证明纤维大小分类定理（仅 $2,3,4$ 三类）。
  \item (C2) 定义仅由“展开（逐位叠加）/折叠（局域重写归一化 + 一次模回卷）”组成的闭合运算 $\oplus$，
  并证明其微观同构于 $(\ZZ_{64},+)$。
  \item (C3) 将 $\oplus$ 编译为一个半径 $2$、相位周期 $90$ 的局域重写 CA（配套规则表与解释器），并对 $64^2$ 输入对穷举验证正确性。
  \item (C4) 构造并可视化一个以 $\oplus$ 为局域相互作用的 $1$D 环形 CA，在微态/空间态/纤维三层同时呈现其时空结构。
\end{itemize}

